\section{A tool call is a typed state transition}\label{sec:C18:contract}
A tool exposes a capability through an interface that an agent can invoke. Its scientific description should include input and output types, preconditions, side effects, failure states, authority requirements, and version identity. A natural-language name and argument schema provide only part of that contract. Two tools with identical JSON shapes can have different resource scopes or effect semantics. A model can produce a syntactically valid call that is operationally wrong.

For environment state space $\mathcal X$, argument space $\mathcal U_j$, observation space $\mathcal O_j$ and receipt space $\mathcal E_j$, represent tool $j$ by a transition kernel
\begin{equation}
 \mathcal T_j(dx',do,d\eta\mid x,u,A),
 \quad u\in\mathcal U_j.
\end{equation}
Here $A$ is current trusted authority, $o\in\mathcal O_j$ an observation and $\eta\in\mathcal E_j$ a typed receipt or uncertainty state. Rejection, timeout and partial completion need explicit outcomes; a partially observed kernel does not give the model access to $x$ or to privileged authorization state. A successful HTTP response is not necessarily proof that the intended business effect occurred; a timeout is not proof that it did not. Chapter~\ref{ch:C13} analyzes the resulting recovery ambiguity.

The model's tool-use decision has at least three stages: select a capability, construct its arguments, and decide how to interpret its result. Toolformer and ReAct are foundational examples of learning or prompting tool-augmented behavior~\cite{schick2023toolformer,yao2022react}. They should not be treated as descriptions of every current runtime. Contemporary systems add discovery, typed schemas, external enforcement, code-based composition, and durable execution. Each addition changes the effective policy and its failure surface.

A contract should identify which fields are exact. Resource identifiers, operation keys, and authority scopes should not be reconstructed from a vague summary when exact values are available. Free-text queries can remain flexible. Output records should separate authoritative data from tool-generated explanations and untrusted external content. Returning a page that contains instructions does not grant those instructions authority over the agent. The wrapper must preserve that distinction when the result enters active context.

For evaluation, define success at the tool boundary and at the task boundary. A schema-valid call can fail a tool precondition. A tool-successful call can be irrelevant to the task. A task can appear solved while the agent produces an unauthorized side effect. These outcomes need separate labels. A single ``tool accuracy'' number is ambiguous unless it specifies whether it measures tool selection, argument validity, execution, or complete justified behavior.

